\phantomsection
\section*{АННОТАЦИЯ}
\addcontentsline{toc}{section}{АННОТАЦИЯ}
\thispagestyle{plain}

Дворянкин~В.~В., разработка моделей и сервиса интерактивного
взаимодействия с обучающимися, магистерская
диссертация, 2026~--- 101~с., 18~рис., 12~табл., 5~листингов,
32~источника, 2~приложения. Руководитель~--- Волков Денис
Андреевич, доцент, к.т.н. Кафедра автоматизированных систем
управления.

Спроектирована архитектура и разработан программный сервис
EduQuiz для интерактивного взаимодействия преподавателя
и обучающихся в синхронных учебных сессиях. Сервис обеспечивает
проведение интерактивных сессий в режиме реального времени
в трёх режимах проведения, индивидуальное повторение учебного
материала по алгоритму SuperMemo SM-2 с автоматическим
формированием оценки качества ответа и психометрический анализ
банка тестовых вопросов средствами классической теории тестов.
Сервис развёрнут в инфраструктуре Российской Федерации
и доступен по доменному имени eduquizz.ru.

Результаты работы могут быть применены в учебном процессе кафедр
российских университетов как программный инструмент интерактивного
тестирования с расширенной геймификацией и алгоритмически
обоснованным повторением учебного материала, дополняющий систему
дистанционного обучения университета.

\newpage
